(1) 两相共存时 $\mu_g(T,P)=\mu_l(T,P)$。微分该等式并用每摩尔基本关系 $d\mu=-s\,dT+v\,dP$，得 $-s_gdT+v_gdP=-s_ldT+v_ldP$，故 $dP/dT=(s_g-s_l)/(v_g-v_l)$，这是严格的Clapeyron关系。
(2) 汽化潜热 $L=T(s_g-s_l)$，所以 $dP/dT=L/[T(v_g-v_l)]$。通常蒸气体积较大且汽化需吸热，因此 $L>0$、$v_g-v_l>0$，饱和压强随温度升高而升高，符合物理观察。
(3) 若 $v_g\simeq RT/P\gg v_l$ 且 $L$ 常数，微分关系化为 $dP/dT=LP/(RT^2)$，即 $d\ln P/dT=L/(RT^2)$。从 $(T_1,P_1)$ 积分到 $(T_2,P_2)$ 得 $\ln(P_2/P_1)=-(L/R)(1/T_2-1/T_1)$，或 $P_2=P_1\exp[(L/R)(1/T_1-1/T_2)]$。
(4) 指数中的 $L/R$ 单位为开尔文，乘以 $1/T$ 后无量纲。临界点附近气液密度差趋零、汽化潜热也趋零，忽略液体体积和恒潜热不再合理；低温区域若蒸气非理想或固相更稳定也须改用相应状态方程。
